% !TEX root = % !TEX root = /Users/nai1ka/Projects/PerfX/Thesis/thesis.tex
\chapter{Conclusion}
\label{chap:conclusion}

This thesis studied performance monitoring and regression detection for Android applications after release. The main problem is that current tools either give no automatic alerts or use fixed thresholds that a developer must set and update by hand for every metric. On a wide range of devices, such absolute thresholds often cause false alerts or missed regressions on some device groups, and a slowdown from a new release can stay unnoticed as long as it remains below the set limit.

To solve this, a complete system was designed and built. A mobile SDK collects performance metrics on real user devices with low overhead and sends them to the backend in batches. A Ktor backend stores the data in ClickHouse and PostgreSQL and assigns a device cohort to each record. A detection pipeline then compares two consecutive releases by the relative P95 shift, grouped by metric, screen, and cohort, and informs developers when a new release makes performance worse. The system was evaluated for runtime overhead and through an end-to-end test of its regression detection. These results, together with the answers to the research questions, are presented in Chapter~\ref{chap:eval}.

The value of this work is that it brings automatic release-to-release regression detection, well established on the server side, to continuous post-release monitoring on real Android devices. Because detection uses the relative change between two releases instead of absolute limits, developers do not need to set or adjust thresholds by hand. The main new element, however, is cohort stratification. Android devices differ widely in capability, so the comparison is separated by device performance cohort, and a regression that appears only on low-end devices is still detected instead of being hidden by faster hardware.